\chapter{Group extensions and the Schur multiplier}
\label{ch:gr02-extensions}

The semidirect product and its internal criterion are established
in \cite[K07.04.semidirect\_product and K07.04.semidirect\_recognition]{d07}.
Here an extension of \(Q\) by \(N\) comes with a specified
inclusion \(N\hookrightarrow E\) and quotient map \(p:E\to Q\).
We identify \(N\) with its image. A splitting is a homomorphism
\(s:Q\to E\) with \(ps=\operatorname{id}_Q\). Its image \(H\) is a
complement, meaning \(NH=E\) and \(N\cap H=\{1\}\). Conversely the
inverse of \(p|_H\) is a splitting for every complement. All groups
in this chapter are discrete, and the abelian kernel \(A\) in a
cohomology group is written additively. Integral group homology
uses the trivial coefficient action unless otherwise specified.

\section{Factor systems and outer actions}

\begin{proposition}[GR02-C04.01. Factor systems for a general kernel]
\label{prop:gr02-factor-system}
Choose a set-theoretic section \(s:Q\to E\) with \(s(1)=1\).
Define \(\alpha_q\in\operatorname{Aut}(N)\) by conjugation by \(s(q)\),
and put \(f(q,r)=s(q)s(r)s(qr)^{-1}\in N\). Then
\begin{align*}
 \alpha_q\alpha_r &=\operatorname{Inn}(f(q,r))\,\alpha_{qr},\\
 f(q,r)f(qr,t)&=\alpha_q(f(r,t))f(q,rt).
\end{align*}
Here \(\operatorname{Inn}(n)(x)=nxn^{-1}\). If \(s'(q)=b(q)s(q)\)
with \(b(1)=1\), the new data are
\[
 \alpha'_q=\operatorname{Inn}(b(q))\alpha_q,\qquad
 f'(q,r)=b(q)\alpha_q(b(r))f(q,r)b(qr)^{-1}.
\]
In particular \(q\mapsto[\alpha_q]\) is a homomorphism
to \(\operatorname{Out}(N)=\operatorname{Aut}(N)/\operatorname{Inn}(N)\)
independent of the section. Conversely, data satisfying the two
identities and the normalizations
\(\alpha_1=\operatorname{id}_N\), \(f(q,1)=f(1,q)=1\)
give an extension on \(N\times Q\) with product
\[
 (n,q)(m,r)=\bigl(n\,\alpha_q(m)f(q,r),qr\bigr).
\]
\end{proposition}

\begin{proof}
Conjugation by \(s(q)s(r)=f(q,r)s(qr)\) gives the first identity.
Expanding \((s(q)s(r))s(t)=s(q)(s(r)s(t))\) and cancelling \(s(qrt)\)
on the right gives the second. Substitution of \(b(q)s(q)\) gives
the displayed change of section. The inner automorphisms form a
normal subgroup of \(\operatorname{Aut}(N)\), since conjugating
\(\operatorname{Inn}(n)\) by any automorphism \(\beta\) gives
\(\operatorname{Inn}(\beta(n))\). Thus the first identity and the
change of section give the outer homomorphism.

For the converse, if \(u\in N\), the two associations of
\((n,q)(m,r)(u,t)\) have first coordinates
\[
 n\alpha_q(m) f(q,r)\alpha_{qr}(u) f(qr,t)
 \quad\text{and}\quad
 n\alpha_q(m)\alpha_q\alpha_r(u)\alpha_q(f(r,t))f(q,rt).
\]
They agree by the two identities, and the second coordinates both
equal \(qrt\). The normalizations give the identity \((1,1)\).
A right inverse to \((n,q)\) can be found with second coordinate
\(q^{-1}\) by solving \(\alpha_q(m)=n^{-1}f(q,q^{-1})^{-1}\).
A left inverse with that second coordinate is obtained by solving
\(m\alpha_{q^{-1}}(n)f(q^{-1},q)=1\).
Associativity makes the two inverses equal. Projection is a
surjective homomorphism with kernel \(N\times\{1\}\).
For the original extension the coordinate map
\((n,q)\mapsto ns(q)\) is bijective and respects the product.
\end{proof}

When \(N=A\) is abelian, \(\alpha_q\) is independent of the lift of
\(q\) and is an ordinary action of \(Q\) on \(A\). An equivalence
between extensions with fixed kernel and quotient is a homomorphism
of middle groups inducing the identity on \(A\) and on \(Q\).
Such a map is injective because its kernel lies in \(A\). For
surjectivity, first lift the quotient coordinate of a target
element and then adjust the lift by an element of \(A\).

\section{Abelian kernels and low-degree cohomology}

We use the inhomogeneous bar differential from
\cite[A02 and A03]{homalg}, with the left \(Q\)-module structure on
\(A\). For \(a:Q\to A\), the two-coboundary is
\[
 (\delta a)(q,r)=a(q)+q\cdot a(r)-a(qr).
\]
A two-cocycle \(c\) satisfies
\[
 c(q,r)+c(qr,t)=q\cdot c(r,t)+c(q,rt).
\]
It is normalized if \(c(1,q)=c(q,1)=0\). The normalized
cocycles give the same \(H^2\) classes as all inhomogeneous
cocycles. Indeed the identity with \(q=1\) shows
\(c(1,r)=c(1,1)\), and with \(t=1\) shows \(c(q,1)=q\cdot c(1,1)\).
Subtract the coboundary of the constant cochain with value
\(c(1,1)\) to normalize. If two normalized cocycles differ by
\(\delta a\), evaluating at \((1,r)\) gives \(a(1)=0\).

\begin{theorem}[GR02-C04.02. Abelian-kernel extensions and \(H^2\)]
\label{thm:gr02-h2-extensions}
Fix a left \(Q\)-action on \(A\). Equivalence classes of extensions
of \(Q\) by \(A\) inducing that action correspond bijectively to
\(H^2(Q,A)\). The zero class corresponds to the split extensions.
\end{theorem}

\begin{proof}
By Proposition~\ref{prop:gr02-factor-system}, a normalized
section gives a normalized cocycle \(c\), and changing the section
adds a coboundary. An extension equivalence preserves the cocycle
of the transported section. Conversely, a normalized cocycle
defines an extension by the product
\[
 (x,q)(y,r)=(x+q\cdot y+c(q,r),qr)
\]
on \(A\times Q\), by the converse in the same proposition.
Its canonical section \(q\mapsto(0,q)\) recovers \(c\), and the
conjugation action on \(A\) is the given action. If
\(c'=c+\delta a\), then \((x,q)\mapsto(x+a(q),q)\) carries the
product for \(c'\) to the product for \(c\) and fixes the two
end groups. Hence the constructions are inverse on classes.

A homomorphic section has cocycle zero. If the class is zero,
choose \(a\) with \(c+\delta a=0\); after the corresponding change
of section its cocycle is zero, so the section is homomorphic.
By the internal criterion this is precisely the semidirect product
for the fixed action.
\end{proof}

\begin{theorem}[GR02-C04.03. Complements and \(H^1\)]
\label{thm:gr02-h1-complements}
For a split extension with abelian kernel \(A\), choose a reference
splitting \(s_0\). The \(A\)-conjugacy classes of complements to
\(A\) in \(E\) are in bijection with \(H^1(Q,A)\), with the class
of \(s_0(Q)\) corresponding to zero.
\end{theorem}

\begin{proof}
Identify \(E=A\rtimes Q\) by \(s_0\). A section is the graph
of a function \(z:Q\to A\), and the semidirect product rule
makes it a homomorphism precisely when
\(z(qr)=z(q)+q\cdot z(r)\); this also implies \(z(1)=0\).
Every complement determines its unique splitting by projection.
Conjugating the graph by \((a,1)\) changes \(z(q)\) to
\(z(q)+a-q\cdot a\). As \(a\) varies these changes are exactly
the one-coboundaries, since they are the negatives of
\(q\mapsto q\cdot a-a\). The description of \(H^1\) in
\cite[A03]{homalg} now gives the bijection.
\end{proof}

The extension correspondence and the choice of reference
splitting are also treated in \cite{conrad-extensions} and
\cite[Theorems 23.3 and 24.3]{lassueur}. The cohomological
Lyndon--Hochschild--Serre sequence is proved in
\cite[A05]{homalg}. The next section constructs the low-degree
homological sequence directly on bar chains.

\begin{example}[GR02-C04.04. Two extensions of \(C_2\) by \(C_2\)]
\label{ex:gr02-c2-extensions}
Give \(A=C_2\) the trivial action of \(Q=C_2=\{1,t\}\).
A normalized cocycle is determined by \(c(t,t)\), which can
take either value in \(C_2\). Indeed, the identity holds whenever
an argument is \(1\), and the triple \((t,t,t)\) gives just
\(c(t,t)=c(t,t)\). Normalized coboundaries have
\((\delta a)(t,t)=2a(t)=0\). The class \(c(t,t)=0\) gives
\(C_2\times C_2\). In the other class the lift \((0,t)\) has
square \((1,1_Q)\), where \(1\) in the first coordinate denotes
the nonzero element of \(A=\mathbb Z/2\mathbb Z\).
It has order four and generates the middle group \(C_4\).
Thus even extensions with trivial action on the kernel can
represent different cohomology classes.
\end{example}

\section{The Hopf formula}

For groups \(U,V\) contained in a group \(E\), the notation
\([U,V]\) denotes the subgroup generated by
\([u,v]=uvu^{-1}v^{-1}\) for \(u\in U\) and \(v\in V\).
In particular \([E,E]\) is the commutator subgroup. The integral
bar complex from \cite[A10]{homalg}, on trivial coefficients,
identifies \(H_1(K;\mathbb Z)\) with the abelianization of \(K\):
in degree one the generators \([k]\) are cycles and the
boundaries are \([l]-[kl]+[k]\). They impose exactly the relation
that \(k\mapsto[k]\) is a homomorphism to an abelian group.

\begin{proposition}[GR02-C04.05. The homological five-term sequence]
\label{prop:gr02-homological-five-term}
For an extension \(1\to N\to E\xrightarrow{p}Q\to 1\), there is
a natural exact sequence
\[
 H_2(E;\mathbb Z)\longrightarrow H_2(Q;\mathbb Z)
 \xrightarrow{d^2_{2,0}}N/[E,N]\longrightarrow
 H_1(E;\mathbb Z)\longrightarrow H_1(Q;\mathbb Z)\longrightarrow 0.
\]
The maps between the homologies of \(E\) and \(Q\) are induced by
\(p\). The map from \(N/[E,N]\) is induced by inclusion.
\end{proposition}

\begin{proof}
Let \(C_\bullet(G)\) denote the integral normalized bar chains
with trivial coefficients, with any tuple containing \(1\) taken
as zero. The unnormalized integral bar chains compute group homology
by \cite[A10]{homalg}. To check normalization in the degrees at issue,
let \(D_\bullet\) be the subcomplex spanned by tuples containing \(1\).
It is a subcomplex because the two faces removing a chosen identity
cancel, and the other faces still contain it. Here \(D_0=0\) and
\(d[1\mid1]=[1]\) kills \(H_1(D)\). A generator of \(D_2\) has
boundary \([1]\). Modulo boundaries from \(D_3\) all generators of
\(D_2\) equal \([1\mid1]\), since
\[
 d[1\mid1\mid g]=[1\mid g]-[1\mid1],
 \qquad
 d[g\mid1\mid1]=[1\mid1]-[g\mid1].
\]
A cycle in \(D_2\) has coefficient sum zero, so \(H_2(D)=0\).
The quotient by \(D_\bullet\) thus has the same homology through
degree two. Projection is surjective on normalized chains in every degree. Put
\(J_\bullet=\ker(C_\bullet(E)\to C_\bullet(Q))\). Since \(J_0=0\),
its first homology is \(J_1/dJ_2\).

The group \(J_1\) is generated by differences \([nx]-[x]\), where
\(n\in N\) and \(x\in E\). Indeed, choose a representative \(s(q)\)
for each fiber, with \(s(1)=1\). Then \([e]-[s(p(e))]\) is a basis
of \(J_1\) as \(e\) ranges over the elements other than the chosen
representatives. Sending this basis element to the class of
\(e\,s(p(e))^{-1}\) defines \(\rho:J_1\to N/[E,N]\), written
additively in the target. It sends every \([nx]-[x]\) to the class
of \(n\), so it is independent of the representatives. Likewise
\(J_2\) is generated by \([na\mid mb]-[a\mid b]\), with \(n,m\in N\):
choose a reference lift for each nonzero base pair and a degenerate
reference lift for a pair mapping to zero. On applying \(d\), the
three differences have \(\rho\)-values equal to the class of \(m\),
minus the class of \(n(ama^{-1})\), and the class of \(n\),
respectively. Their sum is zero.

Conversely, the boundary of \([n\mid x]-[1\mid x]\) identifies
\([nx]-[x]\) with \([n]\) in \(H_1(J)\). The boundary of
\([n\mid m]-[1\mid1]\) gives \([nm]=[n]+[m]\). The boundary of
\([x\mid n]-[x\mid1]\) gives \([xnx^{-1}]=[n]\) by the first identity.
Thus \(n\mapsto[n]\) descends to \(N/[E,N]\) and is inverse to
\(\rho\). Both maps are natural by their formulas on differences.
The homology sequence of
\(0\to J_\bullet\to C_\bullet(E)\to C_\bullet(Q)\to0\) yields the
asserted sequence. The map from \(H_1(J)\) is induced by inclusion,
and the maps from \(H_i(E)\) to \(H_i(Q)\) are induced by \(p\).

To identify the connecting map with the transgression, filter the
bar complex of \(E\) by the number of entries projecting
nontrivially to \(Q\). In filtration degree \(i\), the bottom row
on the first page is \(C_i(Q)\): a tuple of length \(i+1\) with
one entry in \(N\) has filtration-preserving boundary the
difference of two lifts of the same base tuple. These differences
identify every pair of lifts. Indeed, a multiplication by \(N\)
can be moved through successive adjacent positions by inserting
the multiplier in \(N\), and then removed at the end of the
tuple. The bottom-row differential therefore projects to the bar
differential of \(Q\).

In filtration degree zero, the first-page term in total degree one
is \(H_1(N)\). The incoming first differential comes from
chains with precisely one nontrivial projected entry in degree
two and boundary in filtration zero, modulo boundaries.
The chain \([x\mid n]-[xnx^{-1}\mid x]\) has boundary
\([n]-[xnx^{-1}]\) in filtration zero, so these classes impose
conjugation. Conversely any such chain is in \(J_2\), and the
calculation with \(\rho\) shows that its boundary in filtration
zero maps to zero in \(N/[E,N]\). Thus the second-page target
is exactly \(N/[E,N]\).

For a base two-cycle \(z\), choose any lift \(\widetilde z\).
Write \(d\widetilde z\) as a sum of differences \([nx]-[x]\) and
add the chains \([n\mid x]-[1\mid x]\) with the same coefficients
to \(\widetilde z\). The new boundary is in filtration zero and
is the corresponding sum of \([n]\)'s. Its class is both the
connecting map of the short exact sequence and the second
differential \(d^2_{2,0}\) of the filtered bar complex. This fixes
the transgression sign for \(d[a\mid b]=[b]-[ab]+[a]\).
For the spectral-sequence presentation, see \cite[Section 6.8.3]{weibel}.
\end{proof}

\begin{theorem}[GR02-C04.06. Hopf's formula for the Schur multiplier]
\label{thm:gr02-hopf-formula}
For a free presentation \(1\to R\to F\to G\to 1\), the Schur
multiplier \(H_2(G;\mathbb Z)\) identifies naturally with
\[
       (R\cap[F,F])/[F,R].
\]
The natural identification compares different free presentations.
\end{theorem}

\begin{proof}
Fix a free basis \(S\) of \(F\). Its Cayley graph with edges
\(g\to gs\) for \(s\in S\) is a tree by the reduced-word normal
form \cite[K07.01.free\_group\_universal\_property]{d07}.
The augmented chain sequence
over \(\mathbb Z[F]\) has the form
\[
 0\to \bigoplus_{s\in S}\mathbb Z[F]\to\mathbb Z[F]\to\mathbb Z\to0 .
\]
Its boundary sends the edge \(g\to gs\) to \(gs-g\).
Every \(g-1\) is a path boundary, giving exactness at vertices.
A nonzero edge chain has finite support, and in the forest on
those edges a leaf has nonzero boundary coefficient. This proves
injectivity. Thus the sequence is a free resolution of the
trivial left coefficient module, and the Tor description
\cite[A10]{homalg} gives \(H_2(F;\mathbb Z)=0\).

Proposition~\ref{prop:gr02-homological-five-term} identifies
\(H_2(G;\mathbb Z)\) with the kernel of \(R/[F,R]\to F/[F,F]\),
which is the displayed quotient. For a map of free presentations,
the exact sequence is natural. A map between the quotient groups
can be lifted by choosing images of free generators. If two lifts
are chosen, their values differ elementwise by elements of the
target relation subgroup. Those differences are central modulo
commutators with the target free group, so both maps have the
same values on commutators in that quotient. Hence the induced map
on the displayed kernel is independent of the lift. Applied to
the identity of \(G\), naturality gives the canonical comparison
between presentations. See \cite[Theorem 6.8.8]{weibel}.
\end{proof}

\begin{example}[GR02-C04.07. Multipliers of cyclic groups and their products]
\label{ex:gr02-multiplier-cyclic-product}
For \(C_m\), take \(F\) free on one generator. As \([F,F]=1\),
Theorem~\ref{thm:gr02-hopf-formula} gives \(H_2(C_m;\mathbb Z)=0\).
For \(C_m\times C_n\) with \(m,n\geq 1\), present it with free
generators \(x,y\) and relations \(x^m,y^n,[x,y]\). Let \(R\)
be the normal closure of the relations. The quotient
\(D=F/[F,R]\) has central image of \(R\), which contains
\([D,D]\) because \(F/R\) is abelian. Its commutator subgroup is
generated by \(z=[x,y]\) in \(D\). In fact \(F'\) is normally
generated by \([x,y]\), since adjoining that relation makes
the generators commute, and its image in \(D\) is central.
The relations \(x^m,y^n\) have central image, so
\(z^m=z^n=1\), by the commutator rules for central commutators.
For example \([ab,y]=a[b,y]a^{-1}[a,y]=[b,y][a,y]\) in \(D\).

Put \(d=\gcd(m,n)\). If \(d>1\), map \(x,y\) to \(I+E_{12}\)
and \(I+E_{23}\) in the upper unitriangular \(3\times3\) matrices
over \(\mathbb Z/d\mathbb Z\), where \(E_{ij}\) is a matrix unit.
Their \(m\)-th and \(n\)-th powers, respectively, are identities.
Their commutator is \(I+E_{13}\), which is central in this matrix
group and has order \(d\). The image of \(R\) is central because
its normal generators have central images, so the map kills
\([F,R]\). Hence \(z\) has order at least \(d\), and its order
divides \(d\). When \(d=1\) the latter already gives \(z=1\).
The Hopf formula, with \(F'\subseteq R\), now gives
\[
          H_2(C_m\times C_n;\mathbb Z)\cong C_{\gcd(m,n)} .
\]
\end{example}

\section{The universal central extension}

A central extension of \(G\) is an extension with kernel
in the center of its middle group. It is universal if it
admits a unique homomorphism over \(G\) to every central
extension of \(G\). A group is perfect if it equals its
commutator subgroup.

\begin{theorem}[GR02-C04.08. Universal central extensions of perfect groups]
\label{thm:gr02-universal-central-extension}
A group \(G\) has a universal central extension if and only if
it is perfect. When \(G\) is perfect and \(F/R\cong G\) is any
free presentation, the map
\[
       [F,F]/[F,R]\longrightarrow G
\]
is the universal central extension. Its source is perfect,
and its kernel is \(H_2(G;\mathbb Z)\).
\end{theorem}

\begin{proof}
Suppose \(G\) is perfect, and write \(D=F/[F,R]\) and
\(\overline R\) for the image of \(R\). That image is central.
Since \([F,F]\) maps onto \(G\), the subgroup
\(U=[F,F]/[F,R]=[D,D]\) maps onto \(G\), with central kernel
\((R\cap[F,F])/[F,R]\). Theorem~\ref{thm:gr02-hopf-formula}
identifies this kernel. In addition \(D=U\overline R\), so
every commutator of \(D\) is a commutator of elements of \(U\).
It follows that \(U=[D,D]=[U,U]\).

Given a central extension \(E\to G\), lift the free generators
of \(F\) to \(E\), obtaining a homomorphism \(F\to E\) over
\(G\). Its images of the elements of \(R\) lie in the central
kernel, so the map kills \([F,R]\). Restriction yields a
homomorphism \(U\to E\) over \(G\). If two such homomorphisms
are given, their pointwise quotient lies in the central kernel
and is a homomorphism into that abelian group. It must be
trivial because \(U\) is perfect. This proves universality.

Conversely, if \(U\to G\) is universal, the split central
extension \(U/[U,U]\times G\to G\) admits two maps over \(G\)
from \(U\), by using either the zero map or the abelianization
map in the first coordinate. Uniqueness implies \(U=[U,U]\).
Surjectivity then forces \(G\) to be perfect.
See \cite[Theorem 6.9.5]{weibel}.
\end{proof}

For the stable elementary group \(E(R)\), the application in
\cite[B12]{ktheory} identifies \(K_2(R)\) with the corresponding
integral second homology.
